% ex_anexo_iag.tex
% Anexo: Declaração de Uso de Inteligência Artificial Generativa (PRPG/Unicamp)

\chapter*{Anexo 1: Declaração de Uso de Inteligência Artificial Generativa}
\addcontentsline{toc}{chapter}{Anexo 1: Declaração de Uso de Inteligência Artificial Generativa}
\thispagestyle{empty}

% --- Cabeçalho institucional: logo à esquerda, texto à direita ---
\noindent
\begin{minipage}[c]{0.30\textwidth}
  \includegraphics[width=\linewidth, keepaspectratio=true]{figuras/logo_anexo1}
\end{minipage}%
\hfill
\begin{minipage}[c]{0.60\textwidth}
  \raggedleft\small
  Universidade Estadual de Campinas\\
  Pró-Reitoria de Pós-Graduação\\
  \href{https://www.prpg.unicamp.br}{www.prpg.unicamp.br}
\end{minipage}

\vspace{1.5cm}

\begin{center}
  \textbf{DECLARAÇÃO DE USO DE INTELIGÊNCIA ARTIFICIAL GENERATIVA}
\end{center}

\vspace{0.8cm}

Eu, \textbf{Juliane Cristina Helanski Cardoso}, autora do trabalho intitulado \textit{``Tecnografias de um centro de inteligência artificial: seguindo cientistas e engenheiros universidade afora''}, apresentado ao Programa de Pós-Graduação em Ciências Sociais da Universidade Estadual de Campinas (Unicamp), declaro, para os devidos fins de direito e de integridade acadêmica, ter utilizado ferramentas de inteligência artificial generativa (IAG) de forma estritamente auxiliar no desenvolvimento da presente pesquisa e/ou na redação deste trabalho.

As seguintes atividades específicas foram realizadas com o apoio das ferramentas de IAG, sob plena supervisão humana, validação crítica e responsabilidade da autora:

\begin{itemize}
  \item[{[\texttimes]}] Revisão gramatical e edição textual;
  \item[{[\texttimes]}] Auxílio na descrição, estruturação ou execução de procedimentos estatísticos;
  \item[{[\texttimes]}] Suporte na geração ou otimização de códigos de programação;
  \item[{[\texttimes]}] Outras atividades: interlocução argumentativa; construção de diagramas; organização e cruzamento de fontes.
\end{itemize}

As ferramentas de Inteligência Artificial Generativa utilizadas foram:

\begin{enumerate}
  \item \textbf{Claude} (versões Sonnet e Opus, da Anthropic): utilizada como interlocutor argumentativo, para revisão gramatical, edição e organização textual.
  \item \textbf{Claude Code} (versões Sonnet e Opus, da Anthropic): auxílio na descrição, estruturação e execução de procedimentos estatísticos, como análise de rede textual e análise de dados bibliográficos, e na geração ou otimização de códigos de programação (Python e CSS) na análise de áudio e no site da tese.
  \item \textbf{NotebookLM} (Google): organização de conjuntos temáticos de obras e busca do desenvolvimento de ideias na bibliografia.
\end{enumerate}

Em conformidade com a Deliberação CONSU-A-005/2026, declaro estar ciente de que a responsabilidade pelo conteúdo final, veracidade das informações e conclusões deste trabalho é integralmente minha. As ferramentas de inteligência artificial generativa utilizadas não são reconhecidas como coautoras e não possuem qualquer responsabilidade sobre os resultados, dados e conteúdos apresentados.

\vspace{1.5cm}

\noindent Campinas, \today.

\vspace{2cm}

\begin{center}
  \begin{minipage}{0.45\textwidth}
    \centering
    \includegraphics[height=1.2cm, keepaspectratio=true]{figuras/assinatura_juliane}\\[-0.35cm]
    \rule{\textwidth}{0.5pt}\\
    Juliane Cristina Helanski Cardoso\\
    \small Autora
  \end{minipage}
  \hfill
  \begin{minipage}{0.45\textwidth}
    \centering
    \includegraphics[height=1.2cm, keepaspectratio=true]{figuras/assinatura_suely.jpg}\\[-0.35cm]
    \rule{\textwidth}{0.5pt}\\
    Maria Suely Kofes\\
    \small Orientadora
  \end{minipage}
\end{center}